\documentclass[10pt]{resume} % Use the custom resume.cls style
\usepackage{graphicx}
\usepackage{wrapfig}
\usepackage{caption}
%\usepackage{xurl}
\usepackage{breakurl}
\usepackage[colorlinks]{hyperref}
\usepackage{comment}
\usepackage[left=0.3in,top=0.3in,right=0.3in,bottom=0.3in]{geometry}
%\usepackage[hyphenbreaks]{breakurl}
%\usepackage[hyphens]{url}
\newcommand{\tab}[1]{\hspace{.2667\textwidth}\rlap{#1}}
\newcommand{\itab}[1]{\hspace{0em}\rlap{#1}}
\cv{\url{https://faculty.utrgv.edu/sergei.chuprov/}}
\cvs{\href{https://scholar.google.com/citations?hl=en&authuser=1&user=IpuSnGwAAAAJ}{Google Scholar}~~~~~
\href{https://linkedin.com/in/sergei-churpov-889053174}{LinkedIn}
}

\name{Sergei Chuprov}
% Your name
\address{PhD in Computing and Information Sciences}
\address{The University of Texas Rio Grande Valley (UTRGV), Edinburg, TX}
% Your address
%\address{123 Pleasant Lane \\ City, State 12345} % Your secondary addess (optional)

\address{\textbf{Email:} sergei.chuprov@utrgv.edu} % Your phone number and email

\begin{document}

\vspace{-0.3cm}
%----------------------------------------------------------------------------------------
%	EDUCATION SECTION
%----------------------------------------------------------------------------------------
%%%
\begin{comment}
\begin{rSection}{Manifesto}

I am a Ph.D. student in Computer and Information Sciences at the Department of Computer Science at Rochester Institute of Technology. My current research is focused on a security assurance in distributed and multi-agent systems, and ML- and AI-based systems security. 

My comprehensive academic background and valuable research expertise are evidenced by numerous publications and grant awards (please, see below). In 2021, I was awarded the Scholarship of the Russian Federation President for Students Studying Abroad. This Scholarship is designated to encourage students with excellent academic and research achievements for advanced research and cultural exchange experience. I underwent through multiple candidacy selection filters on various levels, from the University to a Russian Federation-wide. My achievements, gained through hard work and with the help of sincere enthusiasm, facilitated to the selection of my candidacy from thousands of others to award the Scholarship. This Scholarship was essential for me, especially in my first PhD year in the US. It helped me to cover my living expenses, which became really high in the end of 2021 due to various economical factors. This financial aid has helped me to focus on my research and academic goals, and to produce numerous research results that will contribute to my PhD thesis. These results include provisional patent application (see Patents section) submitted in August 2022. I highly appreciate the chance to use this Scholarship for the benefit of both myself and the research community.  

% My research is important for industry, and Russian Government made wrong decision. I hope META needs this better.  

Unfortunately, due to the current circumstances occurred in the World since February 2022, The Ministry of Higher Education of the Russian Federation has canceled all the Scholarship award payments to the students studying abroad in the countries, recognized by the Russian Government as ``unfriendly'' on the official level (\href{https://en.wikipedia.org/wiki/Unfriendly_Countries_List}{\nolinkurl{URL}}). I want to express my strongest disapproval with this decision and feel deep misunderstanding. I strongly believe that my research can be beneficial for the industry by making distributed systems more robust, secure, and reliable. I think the Russian Government made the wrongest possible decision by arbitrarily withdraw the intentions to support my research by the reason I chose the ``unfriendly'' country. I believe that Meta's products, services, and end-users will definitely benefit from utilizing my research results. This fellowship will be extremely helpful for me too, as it will give me a chance to gain industry-leading experience and to utilize it in order to pursue my research and academic goals.          
\end{rSection}
\end{comment}

\begin{rSection}{Research Areas of Interest}
      Artificial Intelligence and Machine Learning Systems, Robustness of Intelligent Systems and their Practical Applications, Cybersecurity, Data Quality and Security Evaluation and Analytics
\end{rSection}

\begin{rSection}{Highlights}
\begin{itemize}
\itemsep-0.5em
    %\item Watch my \href{https://youtu.be/3gb9IL-j1NQ?si=YljHl51v5LOUie-b}{short self-presentation video on YouTube}.
    \item 53 peer-reviewed publications, 2 pending patent applications, Google Scholar \textit{\textbf{h}}-index: 9.
    \item My publication list includes renowned peer-reviewed conferences and journals, such as IEEE TDSC journal, IEEE ICC, IEEE Systems journal, IEEE CEC, IEEE WCCI, IEEE INFOCOM (CNERT Workshop), and others.
    %\item My research has been supported by NSF, DoD (USMA), CRDF Global, and other funding agencies.
    \item I have a valuable experience participating in R\&D projects supported by NSF, DoD (USMA), CRDF Global, and other funding agencies.
    \item I received multiple grants and scholarships awarded for the excellence demonstrated in research and education. 
    \item I have 7+ years of Graduate Research and Teaching Assistant experience, mentoring and advising students in various institutions and countries.  
\end{itemize} 
\end{rSection}

\begin{rSection}{Education}
PhD in Computing and Information Sciences \\
{\bf Rochester Institute of Technology (RIT), NY, USA} \hfill \small{\em August 2021 - May 2024}\normalsize\\ 
Department of Computer Science and Department of Cybersecurity\hfill \small\textbf{GPA:} 3.83/4.0\normalsize\\
\textbf{Thesis:} \href{https://www.proquest.com/openview/e54b39dac756a628c82e402e71182a0d/1?pq-origsite=gscholar&cbl=18750&diss=y}{Robust Machine Learning Under Vulnerable Cyberinfrastructure and Varying Data Quality}\\

\begin{comment}
\textbf{Courses completed:}
\scriptsize
\begin{enumerate}
\itemsep-0.8em 
    \item Research Foundations
    \item Quantitative Foundations
    \item Cyberinfrastructure Foundations
    \item Deep Learning
    \item Research Methods in HCI
    \item Systems \& SCADA Security
    \item Foundations of Intelligent Security Systems
    \item Information Security Risk Management
    \item Colloquium on Computational \& Information Sciences
\end{enumerate}
\normalsize
\end{comment}

%PhD in Information Security (\textit{transferred to RIT}) \\
%{\bf Saint-Petersburg National Research University of} \hfill {\em September 2019 - August 2021} \\ \textbf{Information Technologies, Mechanics and Optics, Russia} \hfill \textbf{GPA:} 4.0/4.0 \small \normalsize \\ %\hfill \textit{transferred to RIT} \\
%Faculty of Secure Information Technologies \\

\begin{comment}
\textbf{Courses completed:}
\scriptsize
\begin{enumerate}
\itemsep-0.8em 
    \item English Language
    \item Philosophy and History of Science
    \item Psychology and Teaching in Higher Education
    \item Theory and Methodology in Scientific Research
    \item Patents and Intellectual Property Assurance
    \item Modern Methods in Complex Systems Research  
    \item Teaching Practice
    \item Academic Research
\end{enumerate}
\normalsize
\end{comment}
%PhD in Computer Science (in progress) \\
%{\bf Saint-Petersburg National Research University of} \hfill {\em September 2019 - Present}\\ \textbf{Information Technologies, Mechanics and Optics, Russia} \hfill \textbf{GPA:} 4.0/4.0 (current)\\
%Faculty of Secure Information Technologies\\

\vspace{-0.3cm}
Master's in Information Security \\
{\bf Saint-Petersburg National Research University of} \hfill \small{\em September 2017 - June 2019}\normalsize\\ \textbf{Information Technologies, Mechanics and Optics (ITMO), Russia} \hfill \small\textbf{GPA:} 4.0/4.0 (with honors)\normalsize   
\\ Faculty of Secure Information Technologies\\
\textbf{Thesis:} Assurance of Information in a Mobile Robotic System with Decentralized Control\\
\begin{comment}
\textbf{Courses completed:}
\scriptsize
\begin{enumerate}
\itemsep-0.8em 
\item Information Security Basics
\item Information Security Management
\item Organizational and Legal Mechanisms for Information Security Assurance
\item Foreign Language in Professional Activity (English)
\item Emotional Regulation Status in Professional Activity
\item Work Practice, Research Work
\item Information Security Risks Management
\item Organization of Information Security Service
\item Mathematical Modeling of Objects and Systems
\item Project Management
\item Information Security Technology
\item Technical Means against Industrial Spying
\item Software and Hardware Information Security Means
\item Cryptosystems Modeling
\item Comprehensive Approach to Information Security Assurance
\end{enumerate}
\normalsize
\end{comment}

%$~$\\
\vspace{-0.3cm}
Bachelor's in Information Security \\
{\bf ITMO University, Saint Petersburg, Russia} \hfill \small{\em September 2013 - June 2017}\normalsize \\ 
Faculty of Secure Information Technologies \hfill \small\textbf{GPA:} 3.5/4.0 \normalsize\\
\textbf{Thesis:} Vulnerability Analysis and Classification in Local Networks Based on Microsoft OS Family\\
\begin{comment}
\textbf{Courses completed:}
\scriptsize
\begin{enumerate}
\itemsep-0.8em 
    \item History
    \item Philosophy 
    \item Foreign Language (English)
    \item Economics
    \item Law
    \item Fundamentals of Management Activity
    \item Mathematics
    \item Probability Theory and Mathematical Statistics
    \item Physics
    \item Computer Science
    \item Information Theory
    \item Fundamentals of Information Security
    \item Computer Hardware
    \item Programming Languages
    \item Technology and Methods of Programming
    \item Fundamentals of Electrical Engineering
    \item Electronics and Circuit Design
    \item Data Transmission Networks and Systems
    \item Records Management
    \item Software and Hardware Tools of Information Protection
    \item Technical Information Protection
    \item Cryptographic Methods of Information Protection
    \item Information Technology
    \item Organisational and Legal Support of Information Security
    \item Information Security Management
    \item Health and Safety
    \item Cultural Studies
    \item Social and Political Issues of Modern Society
    \item Ecology
    \item Discrete Mathematics
    \item Advanced Branches of Physics
    \item Concepts of Modern Natural Science
    \item Wave Processes in Media
    \item History and Modern Information Protection System in Russia
    \item Theory of Information Security and Methodology of Information Protection
    \item Information Protection Economics
    \item Information Security Service Organisation and Management
    \item Information Protection Systems in Leading Foreign Countries
    \item Complex Information Protection Systems
    \item Technology of Certification of Information Protection Tools
    \item Information Protection in Database-Management Systems
    \item Protection of Information Processes in Computer Systems
    \item Technology of Vulnerability Detection in Corporate Networks
    \item Protection of Personal Data Processing Systems
    \item Network Operation Systems
    \item Rhetoric
    \item Specialist's Reputation and Corporate Culture
    \item System Theory and System Analysis
    \item Optimization Methods
    \item Graph Theory
    \item Information Protection on the Internet
    \item Modeling of Information Protection Systems 
    \item Information Security Tools Programming
    \item Physical and Wellness Education 
    \item Physical and Wellness Education (Elective course) 
    \item Academic Practical Training (practical training for acquiring basic professional skills and knowledge)
    \item Industrial Placement (practical training for acquiring basic professional skills and experience) including pre-graduation practical training (for execution of graduate qualifying paper)
\end{enumerate}
\normalsize
\end{comment}

\end{rSection}

\vspace{-0.3cm}
\begin{rSection}{Academic Work \& Teaching Experience}

\begin{comment}  
\textbf{B. Thomas Golisano College of Computer Science, RIT}\\
\textbf{Rochester, NY, USA}\hfill\small{\em September 2022 - present}\normalsize\\
{\em Teaching Assistant (part-time)}
\\\textbf{Scope:} Computing and Information Sciences
\\\textbf{Responsibilities:} Assisting with the teaching activities and grading practical assignments of graduate and undergraduate students. Courses:
\begin{itemize}
\itemsep -0.5em 
    \item CSCI-735 Foundations of Intelligent Security Systems. Instructor: Dr. Leon Reznik.
    \item CSCI-636/536 Information Retrieval. Instructor: Dr. Richard Zanibbi.
\end{itemize}
\end{comment}

\textbf{Department of Computer Science, UTRGV}\\
\textbf{Edinburg, TX, USA}\hfill\small{\em September 2024 - present}\normalsize\\
{\em Assistant Professor (tenure-track)}
\\\textbf{Scope \& Responsibilities:} teach undergraduate and graduate courses, conduct research in AI/ML security and robustness, and contribute to departmental and university service activities. \textbf{Courses (full-time instructor):} 
\begin{itemize}
\itemsep-0.5em 
    \item CSCI 2380 Computer Science II (undergraduate)
    \item CSCI 4341 Foundations of Intelligent Security Systems (undergraduate)
    \item CSCI 6370 Foundations of Intelligent Security Systems (graduate)
\end{itemize}


%research and development in the area of Artificial Intelligence- and Machine Learning-based systems security, Data Quality and Security Assurance, Robustness of Machine Learning-end Systems. %\textbf{Results:} 2 pending patent applications; more than 15 published papers in refereed journals and conference proceedings (see Publication List below).

\textbf{B. Thomas Golisano College of Computer Science, RIT}\\
\textbf{Rochester, NY, USA}\hfill\small{\em August 2021 - July 2024}\normalsize\\
{\em Graduate Research Assistant (part-time)}
\\\textbf{Scope \& Responsibilities:} research and development in the area of Artificial Intelligence- and Machine Learning-based systems security, Data Quality and Security Assurance, Robustness of Machine Learning-end Systems. \textbf{Results:} 2 pending patent applications; more than 15 published papers in refereed journals and conference proceedings (see Publication List below).

\textbf{B. Thomas Golisano College of Computer Science, RIT}\\
\textbf{Rochester, NY, USA}\hfill\small{\em August 2022 - July 2024}\normalsize\\
{\em Graduate Teaching Assistant (part-time)}
\\\textbf{Scope \& Responsibilities:} assisting the course instructor in teaching Computer Science-related courses. \textbf{Results:} gained teaching activities experience and knowledge related to building effective communication with the students; holding office hours and grading assignments for the classes of 20 undergraduate and graduate students; providing constructive feedback and communicating students' questions and concerns related to the course content and assignments; supervised student capstone projects and mentored them. 
%add about project supervision with the number of students
%, including the following activities: assisting in the creation of assessment materials; providing constructive feedback on student work; grading assignments, quizzes, exams, and other assessments; holding regular office hours for students; providing one-on-one or small-group tutoring and academic support; assisting students with course-related concerns or challenges. \\
\textbf{Courses:}
\begin{itemize}
\itemsep-0.5em 
    \item CSCI-735 Foundations of Intelligent Security Systems (instructor: \href{https://www.rit.edu/directory/lrvcs-leon-reznik}{Dr. Leon Reznik})
    \item CSCI-536 Information Retrieval, undergraduate section (instructor: \href{https://www.rit.edu/directory/rxzvcs-richard-zanibbi}{Dr. Richard Zanibbi}) 
    \item CSCI-636 Information Retrieval, graduate section (instructor: \href{https://www.rit.edu/directory/rxzvcs-richard-zanibbi}{Dr. Richard Zanibbi})
\end{itemize}

\textbf{Mobile Intelligent Systems Laboratory, ETU ``LETI'' University}
\\
\textbf{Saint Petersburg, Russia}\hfill \small{\em May 2021 - August 2021}\normalsize\\
{\em Researcher (part-time)}
\\\textbf{Scope \& Responsibilities:} research and development in the area of Cyber-Physical Systems security and safety, Intelligent Systems security. Lead a group of 4 developers and engineers; developed models and software tools for Reputation and Trust-based simulation environment to model communication in decentralized IoT networks; applied the developed models and tools on a real-world IoT hardware devices; developed projects' reports and technical documentation. \textbf{Results:} 7 published papers in refereed journals and conference proceedings (see Publication List below).     

%\begin{itemize}
%\itemsep-0.5em 
%    \item \textbf{Teamwork:} the laboratory consists of undergraduate, graduate and postgraduate students.
%\end{itemize}

%\begin{itemize}
%\itemsep-0.5em 
%    \item \textbf{Teamwork:} the laboratory consists of undergraduate, graduate and postgraduate students.
%    \item \textbf{Leadership:} leading a team of 4-6 participants, guiding students in undergraduate studies.
%\end{itemize}

\textbf{Department of Secure Information Technologies, ITMO University}\\
\textbf{Saint Petersburg, Russia}\hfill\small{\em September 2019 - May 2021}\normalsize\\
{\em Teaching Assistant (part-time)}
\\\textbf{Scope \& Responsibilities:} assisting the course instructor in teaching undergraduate and graduate courses in the areas of Information Security and Computer Science. \textbf{Results:} gained experience in teaching laboratory and practical sessions for the classes of 10-20 undergraduate and graduate students; holding office hours and communicating students' questions and concerns; developed and graded assignments for the selected courses; assisted in developing the coursework content for the ``Functional Security and Safety Assurance in Autonomous Vehicles'' Master's degree program, established at ITMO in 2019.
\textbf{Courses:}
\begin{itemize}
\itemsep-0.5em 
    \item Foundations of Information Technology (instructor: Dr. Ilia Viksnin)
    \item Introduction to Imitation Modeling (instructor: Dr. Igor Komarov)
    \item Information Security Technology and Methodology (instructor: \href{https://www.researchgate.net/profile/Viktoria-Korzhuk}{Dr. Victoria Korzhuk})
     \item Foundations of Autonomous Vehicles Safety (instructor: Dr. Ilia Viksnin) 
    \item Autonomous Vehicles Imitation Modeling (instructor: \href{http://geognosis.ifmo.ru/cyril-pshenichny-curriculum-vitae/}{Dr. Cyril Pshenichny})
\end{itemize}

\textbf{Department of Secure Information Technologies, ITMO University}\\
\textbf{Saint Petersburg, Russia}\hfill\small{\em May 2018 - May 2021}\normalsize
\\{\em Researcher (part-time)}
\\\textbf{Scope \& Responsibilities:} research and development in the area of Information Security, Multi-Agent Mobile Robotic Systems Security, Knowledge Representation in Intelligent Systems. \textbf{Results:} developed security and safety models and methods for decentralized and multi-agent mobile robotic systems (see Projects below); developed security-inspired smart factory models in AnyLogic simulation environment; 14 publications in refereed journals and conference proceedings (see Publication List below). 
\end{rSection}

%--------------------------------------------------------------------------------
%    Projects And Seminars
%-----------------------------------------------------------------------------------------------

%\begin{rSection}{Non-academic Work Experience}

%\textbf{OOO Coffee Sirena (Starbucks), Saint-Petersburg}\hfill{\em July 2013 - May 2018}\\{\em Assistant Store Manager (part-time)}
%\\\textbf{Scope:} Catering
%\\\textbf{Job responsibilities:} Responsibility for the economic performance of a coffee shop, staff recruitment and training, leadership of a team of 6-10 persons, resource allocation.
%\\\textbf{Website:} \url{http://www.starbucks.ru/}
%\end{rSection}

\begin{rSection}{Research \& Development Projects}

%{\bf Project Related to Underwater Research} (participant) \\
%\hspace*{\fill} {\em August 2023 - Present}

{\bf Collaborative research: IDEAS lab: ETAUS: Smarter Microbial Observatories for Realtime ExperimentS (SMORES)} (non PI or Co-PI) \\
\hspace*{\fill} {\em August 2023 - Present}
\vspace{-0.3cm}
\begin{itemize}
\itemsep-0.5em 
\item \textbf{Objective:} this recently started project involves the collaboration between RIT, Harvard University, Florida International University, and the University of Georgia to study marine sediments and their role in natural carbon sequestration, using AI and ML to optimize data collection and predictions during field experiments. The project is supported by NSF award \#2321652.  
\end{itemize}

{\bf Enhancing Security and Privacy of the Conventional Federated Learning} (non PI or Co-PI) \\
\hspace*{\fill} {\em September 2022 - Present}
\vspace{-0.3cm}
\begin{itemize}
\itemsep-0.5em 
\item \textbf{Objective:} to develop methods and software tools aimed at detecting compromised and failed local units in the Federated Learning and excluding them from the aggregation. %\textbf{Results:} 4 published papers [C1,C2,C9,C10]; 1 pending patent application [P1]; and 6 poster and other presentations [NP2,NP3,NP5-NP8]. 
The project is supported by the United States Military Academy (USMA) and is accomplished under Grant Number W911NF-20-1-0337. 
\end{itemize}

{\bf Developing Methods and Tools for Machine Learning Robustness Assurance under Vulnerable Cyberinfrastructure and Varying Data Quality} (non PI or Co-PI) \hspace*{\fill} {\em May 2022 - April 2024}
\vspace{-0.4cm}
\begin{itemize}
\itemsep-0.5em 
\item \textbf{Objective:} to develop system integration methods and tools for providing feedback to the Machine Learning system cyberinfrastructure components in order to enhance Machine Learning robustness to input Data Quality variations. 
%\textbf{Results:} 1 pending patent application [P2]; 8 published papers [C5,C6,C11-C16].
%The objective of this project is to develop methods and tools for enhancing ML Robustness, which will allow to significantly improve security and safety of ML systems integrated with network facilities and other modern computing and cyberinfrastructure. 
%In contrast to other approaches, we consider MLINs from the system integration perspective, and propose methods that tide MLIN components together into an interrelated structure with the goal to improve its Robustness to potential data integrity violations. 
The project was supported by the grant from the U.S. Civilian Research \& Development Foundation (CRDF Global) with funding from the United States Department of State. 
\end{itemize}

{\bf Partisan Telegram Application and Operational Security Assessment} (non PI or Co-PI)\\ \hspace*{\fill} {\em February 2022 - July 2022}
\vspace{-0.2cm}
\begin{itemize}
\itemsep-0.8em 
\item \textbf{Objective:} to conduct comprehensive security evaluation of the Partisan Telegram Android application. \textbf{Results:} developed an open-source \href{https://public.opentech.fund/documents/Vulnerability_Assessment_Report_-_OpenTech_Fund_2.pdf}{report} on comprehensive security assessment of the \href{https://github.com/wrwrabbit/Partisan-Telegram-Android}{Partisan Telegram Android OS application}; conducted functional analysis of the application's source code; investigated and verified vulnerabilities that jeopardized user's privacy. The project was supported by the \href{https://www.opentech.fund/news/an-application-and-operational-security-assessment-of-partisan-telegram/}{Open Technology Fund}.    
\end{itemize}

{\bf Enhancing Sensor Network Security with Reputation and Trust-Based Technique} (non PI or Co-PI)\\ \hspace*{\fill} {\em May 2021 - August 2021}
\vspace{-0.2cm}
\begin{itemize}
\itemsep-0.8em 
\item \textbf{Objective:} to develop methods and software tools aimed at detecting malicious nodes in sensor networks using Reputation and Trust indicators. 
%\textbf{Results:} developed software and hardware tools that allows: modeling the communication between real sensor devices; modeling malicious attacks against the quality of the communicated data; verifying the effectiveness of the developed security solutions. 2 published papers [C18,O1]. 
The project was supported by the Ministry of Science and Higher Education of the Russian Federation, \#075-01024-21-02 dated 29.09.2021 (project FSEE-2021-0014)    
\end{itemize}

{\bf Development of Hardware and Software Modules for Critical Objects Monitoring over the Conditions of Uncertainty and Data Incompleteness for Malicious Attacks Prevention and Mitigation}, \textit{Industrial partnership project} (non PI or Co-PI) \hspace*{\fill} {\em September 2019 - October 2020}
\vspace{-0.2cm}
\begin{itemize}
\itemsep-0.5em
\item \textbf{Objective:} to develop software methods and tools implemented on sensor-equipped hardware modules aimed at monitoring and detecting hazardous environmental events in the conditions of malicious attacks. \textbf{Results:} developed intelligent security evaluation models, risk evaluation models and methods, and software and hardware prototypes, implemented by Murmansk Commercial Seaport OAO, Russia. The work was supported by the Ministry of Science and Higher Education of the Russian Federation under the agreement No. 05.605.21.0189. %\href{https://fcpir.ru/participation_in_program/contracts/05.605.21.0189/}{\nolinkurl{No. 05.605.21.0189}}. 
\end{itemize}

{\bf Development of a Quantum Secure Hybrid Platform for Distributed Cyber-Physical Systems Control} (non PI or Co-PI) \hspace*{\fill} {\em May 2018 - August 2019}
\begin{itemize}
\itemsep-0.5em
\item \textbf{Objective:} to develop methods and tools aimed at assuring data security and confidentiality in a group of autonomous mobile robotic devices equipped with quantum key generation facilities. \textbf{Results:} developed Multi-Agent based security models and methods for mobile robotic system equipped with quantum key generation facilities, and mobile robotic devices prototypes, implemented at ITMO University for further research and educational purposes. The work was financially supported by the Government of Russian Federation, Grant 074-U01.
\end{itemize}

{\bf Development of Experimental Testbed for Smart City Control and Security Algorithms Research and Verification} (non PI or Co-PI) \hfill \hfill {\em May 2018 - March 2019}
\begin{itemize}
\itemsep-0.5em
\item \textbf{Objective:} to develop physical testbed for conducting real-world empirical verification of the developed IoT security solutions deployed in decentralized systems, such as Intelligent Transportation Systems. 
%\textbf{Results:} developed physical testbed equipped with IoT communication facilities and 10 autonomous vehicle models for security and optimization algorithms verification. Published 8 papers [J7,J8,C23,C25-C29]. 
The video of the developed testbed can be found in a public access: \href{https://vk.com/video-175152723_456239017}{\nolinkurl{URL}}.
\end{itemize}
\end{rSection}


\begin{comment}
\begin{rSection}{Educational Projects}
\textbf{Development of a cryptosystem model for asymmetric encryption}
\\%A small project on the course of Cryptosystem Modeling.
A model of a cryptosystem on the Python language was developed, allowing to perform secure data transmission with session key and using AES and RSA algorithms.


\textbf{Development of an integrated information security system for an enterprise}
\\A model of an integrated information security system was developed in accordance with the standards and requirements adopted in the Russian Federation.


\textbf{Development of an economically feasible information protection system for an enterprise}
\\An assessment of the damage caused by the realization of possible threats to information security was carried out, a risk assessment was conducted and the appropriateness of the proposed information protection tools was evaluated.
\end{rSection}
\end{comment}



%\begin{rSection}{Other projects}

\newpage


%\begin{rSection}{Grant Applications Writing Experience}

\begin{comment}
\vspace{-1mm}
\textbf{NSF grant for participation in POWDER Mobile and Wireless Week 2023} \hfill {\em 2023} \\
\textit{Supported}, covered travel and participation expenses in the event, \$1,800 

\vspace{-1mm}
\textbf{Amazon AWS Cloud Credit for Research, awarded for one year} \hfill {\em 2022}  \\
\textit{Supported}, covered expenses to use AWS facilities for research, \$4,000

\vspace{-1mm}
\textbf{(2) Saint-Petersburg’s Academic and Industrial Institutions Students’ Grant} \hfill {\em 2022, 2019} \\
\textit{Supported}, covered work expenses, \$1,000 
\end{comment}

\begin{comment}
\vspace{-1mm}
\textbf{Amazon AWS Cloud Credit for Research}, ``Empirical Study of Industrial Machine Learning End Systems Robustness to Data Quality Degradation'', PI \hfill {\em 2022} \\
\textit{Supported}, \$4,000 in AWS credits awarded. \\
\textbf{Summary of contributions:} developed research motivation and problem description, formulated research questions, developed research agenda, developed specification and empirical study design.

\textbf{Ministry of Science and Higher Education of the Russian Federation (\#05.605.21.0189)}, ``Development of Hardware and Software Modules for Critical Objects Monitoring over the Conditions of Uncertainty and Data Incompleteness for Malicious Attacks Prevention and Mitigation'', neither PI or Co-PI (Research Assistant)
\hfill {\em 2019} \\
\textit{Supported}, \$500,000 awarded. \\
\textbf{Summary of contributions:} developed sections related to risk management models, created figures, performed editing.
\end{comment}

%\end{rSection}


%\end{rSection}
%\newpage
\begin{rSection}{Grants and awards}

%\textbf{Paper ``Federated Learning with Trust Evaluation for Industrial Applications'' included into the 10 distinguished conference papers submitted to the IEEE CAI 2023} \\ Santa Clara, CA, USA. \hfill {\em June 2023}


\textbf{\href{https://teachaccess.org/teach-access-fellowship-program/}{2025 Teach Access Fellowship}} \hfill {\em December 2024} \\ USA
  
\vspace{-1mm}
\textbf{IEEE CAI 2023 Student Conference Grant} \hfill {\em June 2023} \\ Santa Clara, CA, USA 

\vspace{-1mm}
\textbf{Gold Medal (1st prize) for the poster presented at the UPSTAT 2023 conference} \hfill {\em April 2023} \\ Rochester, NY, USA 

\vspace{-1mm}
\textbf{NSF grant for participation in POWDER Mobile and Wireless Week 2023} \hfill {\em January 2023} \\
Salt Lake City, UT, USA 

\vspace{-1mm}
\textbf{Amazon AWS Cloud Credit for Research, awarded for one year} \hfill {\em August 2022}  \\
Rochester, NY, USA 

\vspace{-1mm}
\textbf{Saint-Petersburg’s Academic and Industrial Institutions Students’ Grant} \hfill {\em (2): 2019, 2022} \\
Saint-Petersburg, Russia  

\vspace{-1mm}
\textbf{IEEE INFOCOM 2022 Student Conference Grant} \hfill {\em May 2022} \\
London, UK 

\vspace{-1mm}
\textbf{Rochester Institute of Technology PhD Research and Tuition Scholarship} \hfill {\em (3): 2021-2023} \\
Rochester, NY, USA 

\vspace{-1mm}
\textbf{Scholarship of the Russian Federation President for Students of Higher Education} \hfill {\em May 2021} \\ \textbf{Universities Studying Abroad in 2021/2022 Academic Year} \\
Saint-Petersburg, Russia 

\vspace{-1mm}
\textbf{Scholarship of the Russian Federation President and Government for the} \hfill {\em (2): 2021, 2022} \\ \textbf{Students of Priority Education Directions} \\
Saint-Petersburg, Russia 

%\textbf{Scholarship of the Russian Federation President and Government for the}\\ \textbf{Students of Priority Education Directions in 2020/2021 Academic Year} \\
%Saint-Petersburg, Russia  \hfill {\em May 2020} 

%\textbf{Saint-Petersburg's Academic and Industrial Institutions Students' Grant Award} \\
%Saint-Petersburg, Russia \hfill {\em October 2019}

\vspace{-1mm}
\textbf{ITMO University's Best Student's Research MS Thesis Award} \hfill {\em August 2019} \\
Saint-Petersburg, Russia 

\vspace{-1mm}
\textbf{Scholarship of the Russian Federation Government} \hfill {\em September 2019} \\
Saint-Petersburg, Russia 
\end{rSection}


\begin{rSection}{Invited Talks \& Sessions}
    \textbf{Invited Talk (remote) at \href{https://sites.google.com/view/t2p-workshop/home}{From Theory to Practice (T2P)} Worshop, Kigali, Rwanda}\hfill\small{\em 16 April 2024}\normalsize\\
    \textbf{Title:} \textit{Security and Robustness of Machine Learning under Vulnerable Cyberinfrastructure and Varying Data Quality}\\

    \textbf{Invited Session at \href{https://community.amstat.org/rochester/events/upstats-2024}{UPSTAT 2024} Conference, NY, USA}\hfill\small{\em 13 April 2024}\normalsize\\
    \textbf{Title:} \textit{Improving Federated Learning Robustness towards Security Violations and Data Quality Degradations OR How to Learn Better via Extracting Knowledge and Accumulating History?}\\
    Organized together with \href{https://www.rit.edu/directory/lrvcs-leon-reznik}{Dr. Leon Reznik} and Raman Zatsarenko (PhD student, RIT)
\end{rSection}


\begin{rSection}{Service, Volunteering \& Mentoring Experience}

\vspace{-10pt}
\item ACADEMIC SERVICE:

\textbf{IEEE Transactions on Artificial Intelligence}\\
{\em Reviewer: 2025}

\textbf{2025 International Joint Conference on Neural Networks (IJCNN 2025)}\\
{\em Reviewer: 2025}

\textbf{IEEE IoT Journal}\\
{\em Reviewer: 2024, 2025}

\textbf{MDPI Sensors}\\
{\em Reviewer: 2024, 2025}

\textbf{MDPI Network}\\
{\em Reviewer: 2024, 2025}

\textbf{MDPI Electronics}\\
{\em Reviewer: 2024, 2025}

\textbf{Journal of Data and Information Quality}\\
{\em Reviewer: 2024, 2025}

\textbf{IEEE World Congress on Computational Intelligence (WCCI) 2024: The International Joint Conference on Neural Networks (IJCNN)}\\
{\em Reviewer: 2024}

\textbf{The 30th International Conference on Neural Information Processing (ICONIP 2023)}\\
{\em Reviewer: 2023}

\item VOLUNTEERING:

\href{https://www.utrgv.edu/es2/}{\textbf{Engaged Scholar Symposium}}\\
\textbf{The University of Texas Rio Grande Valley, TX, USA}\hfill\small{\em 2025}\normalsize
\\{\em Abstract Reviewer and Judge}
\\ Evaluated and provided feedback for the abstracts and projects presented by students.

\href{https://geniusolympiad.org/}{\textbf{GENIUS Olympiad}}\\
\textbf{Rochester Institute of Technology, NY, USA}\hfill\small{\em (2): 2023, 2024}\normalsize
\\{\em Volunteer Judge}
\\ Evaluated and graded 12 science projects (each year) presented at 2023 and 2024 GENIUS Olympiad. Provided constructive feedback on the projects and presentations to the participants.

\item STUDENT MENTORING:

\textbf{Angel Peredo, PhD}\\
\textbf{The University of Texas Rio Grande Valley, Edinburg, TX, USA}\hfill\small{\em 2025}\normalsize

\textbf{Danny McClain, BS}\\
\textbf{The University of Texas Rio Grande Valley, Edinburg, TX, USA}\hfill\small{\em 2025}\normalsize

\textbf{Ei Ei Nyein Chan, MS}\\
\textbf{The University of Texas Rio Grande Valley, Edinburg, TX, USA}\hfill\small{\em 2025}\normalsize

\textbf{Damian Villarreal, BS}\\
\textbf{The University of Texas Rio Grande Valley, Edinburg, TX, USA}\hfill\small{\em 2025}\normalsize

\textbf{Bryn Ramirez, MS}\\
\textbf{The University of Texas Rio Grande Valley, Edinburg, TX, USA}\hfill\small{\em 2025}\normalsize

\textbf{Nathaniel Deleon, BS}\\
\textbf{The University of Texas Rio Grande Valley, Edinburg, TX, USA}\hfill\small{\em 2025}\normalsize

\textbf{Kayla Luevanos, BS}\\
\textbf{The University of Texas Rio Grande Valley, Edinburg, TX, USA}\hfill\small{\em 2025}\normalsize

\textbf{Andrew Alwarez, BS}\\
\textbf{The University of Texas Rio Grande Valley, Edinburg, TX, USA}\hfill\small{\em 2025}\normalsize

\textbf{Gael Marquez, BS}\\
\textbf{The University of Texas Rio Grande Valley, Edinburg, TX, USA}\hfill\small{\em 2025}\normalsize

\textbf{Adrian Pena, MS}\\
\textbf{The University of Texas Rio Grande Valley, Edinburg, TX, USA}\hfill\small{\em 2024, 2025}\normalsize

\textbf{Pedro Fonseca, PhD}\\
\textbf{The University of Texas Rio Grande Valley, Edinburg, TX, USA}\hfill\small{\em 2024, 2025}\normalsize

\textbf{Eduardo Dominguez, BS}\\
\textbf{The University of Texas Rio Grande Valley, Edinburg, TX, USA}\hfill\small{\em 2024, 2025}\normalsize \\
Supported by UTRGV Engaged Scholar Award for Undergraduate Research (with Dr. Chuprov as a mentor) 

\textbf{Carlos Sepulveda, BS}\\
\textbf{The University of Texas Rio Grande Valley, Edinburg, TX, USA}\hfill\small{\em 2024}\normalsize

\textbf{Harshil Pravinbhai Patel, MS}\\
\textbf{Rochester Institute of Technology, NY, USA}\hfill\small{\em 2024}\normalsize
\\MS capstone project: ``Improving Federated Learning Security with Trust Evaluation to Detect Adversarial Attacks''


\textbf{Shweta Sandip Sharma, MS}\\
\textbf{Rochester Institute of Technology, NY, USA}\hfill\small{\em 2024}\normalsize
\\MS capstone project: ``Analyzing Effects of Adversarial Attacks on Classifier Performance''

\textbf{Gaurav Thakur, MS}\\
\textbf{Rochester Institute of Technology, NY, USA}\hfill\small{\em 2024}\normalsize
\\MS capstone project: ``Resilience of Federated Learning Models in Image Classification Under Data Degradation''

\textbf{Vedika Vishwanath Painjane, MS}\\
\textbf{Rochester Institute of Technology, NY, USA}\hfill\small{\em 2024}\normalsize
\\MS capstone project: ``Analyzing the Effect of Data Poisoning on Semantic Segmentation Using Federated Learning''

\textbf{Anirudh Narayanan, MS}\\
\textbf{Rochester Institute of Technology, NY, USA}\hfill\small{\em 2024}\normalsize
\\MS capstone project: ``Soccer Games Automatic Video Highlights Generator'', [\href{https://github.com/Anirudhrn98/Event_Detection_Soccer}{GitHub}] 

\textbf{Moinuddin Memon, MS}\\
\textbf{Rochester Institute of Technology, NY, USA}\hfill\small{\em 2023}\normalsize
\\MS capstone project: ``Knowledge-Based Client Filtering Algorithm for Federated Learning'' (see [C10]) 

\textbf{Ilya Sakhno, BS}\\
\textbf{ITMO University, Saint Petersburg, Russia}\hfill\small{\em 2021}\normalsize
\\BS qualification project: ``Unified Security Profile for Contactless Payment Technology''

\textbf{Egor Marinenkov, BS}\\
\textbf{ITMO University, Saint Petersburg, Russia}\hfill\small{\em 2020}\normalsize
\\BS qualification project: ``Detecting Data Integrity Violations in Cyber-Physical Systems using Game Theory Approach'' (see [J5] and [C31])

\end{rSection}

\begin{rSection}{Additional Research Experience}
Summer research program in Cybersecurity under the supervision of \href{https://www.rit.edu/directory/lrvcs-leon-reznik}{Dr. Leon Reznik}\\
{\bf RIT, NY, USA} \hfill {\em June 2019 - August 2019}\\ 
\textbf{Results:} 3 published papers [J6,C19,C20] related to security and safety assurance in a group of autonomous vehicles with Reputation, Trust, and Data Quality indicators.\\

\end{rSection}

\begin{comment}
\begin{rSection}{Special Education in English}
{\bf Alpha-Dialogue Intensive Language Session, Pargas, Finland} \hfill {\em 3-12th January 2013}
\\ Intensive English language course. Became the winner of the program for the best progress in the English language in my age category.


{\bf Embassy English Language School, Hastings, UK} \hfill {\em 25th March - 7th April 2012}
\\ Intensive English course during spring school holidays. At the end of the course, received a certificate of B2 level and positive feedback from teachers.

{\bf EF International Language School, Dublin, Ireland} \hfill {\em 4-22th July 2011}
\\ Successfully completed a full-time coourse at EF Summer School of English Language. Received a certificate of A2 level.
\end{rSection}
\end{comment}

\begin{comment}
\begin{rSection}{Additional Training Courses}
``Graduate Teaching Assistant Foundations'' training\\
{\bf RIT, Rochester, NY} \hfill {\em December 2022} \\
Online training included 3 modules that covered: strategies and resources for supporting academic integrity; ways to ensure equitable opportunities for students; digital tools for classroom learning goals; instructions on identifying and responding to microaggressions in the classroom, and other topics.

Online intensive course: ``Blended Learning: Digital Teaching Competencies''\\
{\bf ITMO University, Saint Petersburg, Russia} \hfill {\em 24 - 29 August 2020} \\
72 hours of online intensive lectures, practice, and master-classes with a final project defense. Confirmed by a state-recognized certificate.

In-person intensive course: ``Foundations of Public Speaking in the 21st Century''\\
{\bf ITMO University, Saint Petersburg, Russia} \hfill {\em 13 March - 24 April 2019}\\ 
72 hours of intensive lectures, practice, and master-classes with final public-speaking project defense. Confirmed by a state-recognized certificate.
\end{rSection}
\end{comment}

%\newpage
\begin{comment}
\begin{rSection}{Skills \& Tools}

\begin{tabular}{ @{} >{\bfseries}l @{\hspace{6ex}} l }
Modeling \& Analysis \ & Python Libraries (NumPy, pandas, scikit-learn, etc),  Matlab, \\ & AnyLogic, IBM SPSS, Network Simulator (2,3), \\ & Simulator of Urban Mobility (SUMO), CARLA \\
Deep Learning \& Computer Vision  & Python Frameworks (TensorFlow, PyTorch, Keras, Flower) \\
Software Development & Python, C/C++, Git \\  
%Languages \ & English, Russian (native) \\
\end{tabular}
\end{rSection}
\end{comment}

\newpage

\begin{comment}
\begin{rSection}{Media}

\textbf{Related publications in social media:}
\item ITMO’s R\&D Project Results and Their Applications: \href{https://news.itmo.ru/en/science/cyberphysics/news/9255/#:~:text=100%25%20of%20attacks.%E2%80%9D-,A%20drone%20model%20testing%20stand,-Sergey%20Chuprov%E2%80%99s}{\nolinkurl{URL}}

\end{rSection}

\begin{rSection}{Referees}

\textbf{\href{https://www.rit.edu/directory/lrvcs-leon-reznik}{Dr. Leon Reznik}, advisor,} Professor of Computer Science (primary affiliation) and Computing Security (secondary affiliation) at RIT, NY. \textbf{Email:} leon.reznik@rit.edu 

\textbf{\href{https://www.rit.edu/directory/rxzvcs-richard-zanibbi}{Dr. Richard Zanibbi}}, Professor of Computer Science and Director of the Document and Pattern Recognition Lab at RIT, NY. \textbf{Email:} rxzvcs@rit.edu

\textbf{\href{https://engineering.louisville.edu/faculty/roman-v-yampolskiy/}{Dr. Roman Yampolskiy}}, Associate Professor of Computer Science and Engineering and Director of Cyber Security Laboratory in Speed School of Engineering, University of Louisville, KY. \textbf{Email:} roman.yampolskiy@louisville.edu

\textbf{\href{https://www.ifi.uzh.ch/en/spins/people/oliveira.html}{Dr. Ivan De Oliveira Nunes}}, Assistant Professor of Cybersecurity at the University of Zurich (UZH), Switzerland. \textbf{Email:} nunes@ifi.uzh.ch
    
\end{rSection}
\end{comment}

%~\\
\newpage

\centering \Large \textbf{PUBLICATION LIST}
\normalsize

\begin{rSection}{Patents}
\begin{enumerate}
\itemsep-0.5em

\item Chuprov, S., \& Reznik, L. ``Federated Learning with A Compromised Unit Exclusion from Receiving Global Model Updates''.\\
\scriptsize \textit{Provisional application filed on January 13, 2023}\normalsize

\item Chuprov, S., \& Reznik, L. ``Network Adjustment based on Machine Learning End System Performance Monitoring Feedback''.\\
\scriptsize\textit{Provisional application filed on September 14, 2022}\normalsize

\end{enumerate}
\end{rSection}

\begin{rSection}{Refereed Journal Publications}
\begin{enumerate}
\itemsep-0.5em

\item Berezovskaya, O., Chuprov, S., Neverov, E., \& Sadreev, E. (2023). ``Review and Comparison of Lightweight Modifications of the AES Cipher for a Network of Low-Power Devices'' in Automatic Control and Computer Sciences. 2022, no. 56, pp. 994-1006. doi: 10.3103/S0146411622080028. \href{https://link.springer.com/article/10.3103/S0146411622080028}{\nolinkurl{URL}}

\item Melnikov, T., Chuprov, S., Lazarev, E., Gataullin, R., \& Viksnin, I. (2022). ``Improving Reputation and Trust-Based Approach with Reliability Indicators for Autonomous Vehicles Intergroup Communication'' in Tomsk State University Journal of Control and Computer Science. 2022, no. 61, pp. 108-117.
\href{http://journals.tsu.ru/informatics/en/&journal_page=archive&id=2302#:~:text=Improving%20reputation%20and,vehicles%20intergroup%20communication}{\nolinkurl{URL}}

\item Marinenkov, E., Chuprov, S., Tursukov, N., Kim, I., \& Viksnin, I. (2022). ``Study on Destructive Informational Impact in Unmanned Aerial Vehicles Intergroup Communication'' in Symmetry. 2022. Vol. 14, no. 8, pp. 1-18. \href{https://www.mdpi.com/2073-8994/14/8/1580/htm}{\nolinkurl{URL}} \scriptsize{\textcolor{red}{\textit{Published in Special Issue}}}\normalsize

\item Buzina, E., Chuprov, S., Tursukov, N., Belyaev, P., \& Viksnin, I. (2022). ``Algorithm for Uniform Coverage of the Monitoring Territory by Unmanned Aerial Vehicles'' in LETI Transactions on Electrical Engineering \& Computer Science. 2022. Vol. 15, no. 5/6. P. 41–50. doi: 10.32603/2071-8985-2022-15-5/6-41-50. -- \textit{in Russian.} \href{https://izv.etu.ru/ru/arhive/2022/5-6/41-50}{\nolinkurl{URL}}

\item Berezovskaya, O., Chuprov, S., Neverov, E., \& Sadreev., E. (2022). ``Review and Comparison of AES Lightweight Modifications for a Low-Power Devices Network'' in Problems of Information Security. Computer Systems. no. 2, pp. 35-50. doi: 10.48612/jisp/gp9v‑96dh‑32\\
v1. -- \textit{in Russian.} \href{https://jisp.ru/article/obzor-i-sravnenie-legkovesnyh-modifikatsij-shifra-aes-dlya-seti-malomoshhnyh-ustrojstv/}{\nolinkurl{URL}}

\item Viksnin, I., Marinenkov, E., \& Chuprov, S. (2022). ``A Game Theory approach for communication security and safety assurance in cyber-physical systems with Reputation and Trust-based mechanisms'' in Scientific and Technical Journal of Information Technologies, Mechanics and Optics, vol. 22, no. 1, pp. 47–59. doi: 10.17586/2226-1494-2022-22-1-47-59. \href{https://ntv.ifmo.ru/en/article/20971/primenenie_teorii_igr_dlya_obespecheniya_bezopasnosti_kommunikacii_kiberfizicheskoy_sistemy_s_ispolzovaniem_mehanizmov_reputacii_i_doveriya.htm}{\nolinkurl{URL}}

\item Chuprov, S., Gataullin, R., \& Viksnin, I. (2021). ``Vulnerabilities Assesment in Mobile Robotic Control System'' in Proceedings of the St. Petersburg State Electrotechnical University LETI, vol. 10, pp. 70-75. -- \textit{in Russian}. \href{https://izv.etu.ru/ru/arhive/2021/10/70-75}{\nolinkurl{URL}}

\item Khokhlov, I., Reznik, L., \& Chuprov, S. (2020). ``Framework for Integral Data Quality and Security Evaluation in Smartphones'' in IEEE Systems Journal, vol. 15, no. 2, pp. 2058-2065, doi: 10.1109/JSYST.2020.2985343. \href{https://ieeexplore.ieee.org/abstract/document/9080556}{\nolinkurl{URL}}

\item Chuprov, S., Viksnin, I., Kim, I., Tursukov, N., \& Nedosekin, G. (2020). Empirical Study on Discrete Modeling of Urban Intersection Management System. International Journal of Embedded and Real-Time Communication Systems (IJERTCS), vol. 11(2), pp. 16-38, doi: 10.4018/IJERTCS.2020040102. \href{https://www.igi-global.com/article/empirical-study-on-discrete-modeling-of-urban-intersection-management-system/258082}{\nolinkurl{URL}}

\item Chuprov, S., Viksnin, I., Kim, I., Marinenkov, E., Usova, M., Lazarev, E., Melnikov, T., \& Zakoldaev, D. (2019). Reputation and Trust Approach for Security and Safety Assurance in Intersection Management System. Energies, 12(23), 4527, doi: 10.3390/en12234527. \href{https://www.mdpi.com/1996-1073/12/23/4527}{\nolinkurl{URL}}
\end{enumerate}

\end{rSection}

\begin{rSection}{Refereed Conferences' Proceedings and presentations}
\begin{enumerate}
\itemsep-0.5em

\item Chuprov, S., Reznik, L., \& Grigoryan, G. (2022). ``Study on Network Importance for ML End Application Robustness'' in 2023 IEEE International Conference on Communications. \\
\scriptsize \textit{Accepted for presentation and publication. Will be published after June 2023.} \normalsize

\item Chuprov, S., Satam, A. N., \& Reznik, L. (2022). ``Are ML Image Classifiers Robust to Medical Image Quality Degradation?'' in 2022 IEEE Western New York Image and Signal Processing Workshop (WNYISPW), 2022, pp. 1-4, doi: 10.1109/WNYISPW57858.2022.9983488.
\href{https://ieeexplore.ieee.org/document/9983488}{\nolinkurl{URL}}

\item Chuprov, S., Khokhlov, I., Reznik, L., \& Manghi, K. (2022). ``Multi-Modal Sensor Selection with Genetic Algorithms'' in 2022 IEEE Sensors, 2022, pp. 1-4, doi: 10.1109/SENSORS52175.2022.9967296.
\href{https://ieeexplore.ieee.org/document/9967296}{\nolinkurl{URL}}

\item Khokhlov, I., Chuprov, S., \& Reznik, L. (2022). ``Integrating Security with Accuracy Evaluation in Sensors Fusion'' in 2022 IEEE Sensors, 2022, pp. 1-4, doi: 10.1109/SENSORS52175.2022.9967235.
\href{https://ieeexplore.ieee.org/document/9967235}{\nolinkurl{URL}}

\item Chuprov, S., Khokhlov, I., Reznik, L., \& Shetty, S. (2022). ``Influence of Transfer Learning on Machine Learning Systems Robustness to Data Quality Degradation'' in 2022 International Joint Conference on Neural Networks (IJCNN 2022), 2022, pp. 1-8, doi: 10.1109/IJCNN55064.2022.9892247. \href{https://ieeexplore.ieee.org/abstract/document/9892247}{\nolinkurl{URL}}

\item Chuprov, S., Reznik, L., Obeid, A., \& Shetty, S. (2022). ``How Degrading Network Conditions Influence Machine Learning End Systems Performance?'' in IEEE INFOCOM 2022 - IEEE Conference on Computer Communications Workshops (INFOCOM WKSHPS), 2022, pp. 1-6, doi: 10.1109/INFOCOMWKSHPS54753.2022.9798388. \href{https://ieeexplore.ieee.org/abstract/document/9798388?casa_token=F_FXxDTupigAAAAA:OyMilDynWbzOuAZ4SjcNAbJR9_2PQubP-8I1GWdwoHncOIxMz7boIrtaOs2ik7b6BY_C72m2DQ}{\nolinkurl{URL}}

\item Chuprov, S., Gataullin, R., Neverov, E., Belyaev, P., Kim, I., \& Viksnin, I. (2022). ``Police Office Model Performance and Security Evaluation in a Simulated Group of Mobile Robots'' in The 5th International Conference on Future Networks \& Distributed Systems (ICFNDS 2021). Association for Computing Machinery, New York, NY, USA, pp. 606–615, doi: 10.1145/3508072.3508195. \href{https://dl.acm.org/doi/abs/10.1145/3508072.3508195?casa_token=emU3vhg0J98AAAAA:DwcQ-b1JVuMRRTeGcLDn6ZFtiSzuP4CL88yTMTYZIYGehMFD-kwUEAlGoYH-j-yRxkK1xONmBQo}{\nolinkurl{URL}}

\item Chuprov, S., Viksnin, I., Kim, I., Melnikov, T., Reznik, L., \& Khokhlov, I. (2021). ``Improving Knowledge Based Detection of Soft Attacks Against Autonomous Vehicles with Reputation, Trust and Data Quality Service Models'' in 2021 IEEE International Conference on Smart Data Services (SMDS), pp. 115-120. \href{https://ieeexplore.ieee.org/document/9592379}{\nolinkurl{URL}}

\item Domnitsky, E., Mikhailov, V., Zoloedov, E., Alyukov, D., Chuprov, S., Marinenkov, E., \& Viksnin, I. (2021). ``Software Module for Unmanned Autonomous Vehicle’s On-board Camera Faults Detection and Correction'' in CEUR Workshop Proceedings, Vol. 2893, pp. 1-10. \href{http://ceur-ws.org/Vol-2893/paper_18.pdf}{\nolinkurl{URL}}

\item Lyakhovenko, Y., Viksnin, I., \& Chuprov, S. (2021). ``Integrating Smart Contracts into Smart Factory Elements’ Informational Interaction Model'' in CEUR Workshop Proceedings, Vol. 2893, pp. 1-6. \href{http://ceur-ws.org/Vol-2893/short_4.pdf}{\nolinkurl{URL}}

\item Usova, M., Viksnin, I., \& Chuprov, S. (2021). ``Informational Messages and Space Models Application in Smart Factory Concept'' in CEUR Workshop Proceedings, Vol. 2893, pp. 1-8. \href{http://ceur-ws.org/Vol-2893/short_3.pdf}{\nolinkurl{URL}}

\item Khanh, T.D., Komarov, I., Don, L.D., Iureva, R., \& Chuprov, S. (2020). ``TRA: Effective Authentication Mechanism For Swarms Of Unmanned Aerial Vehicles'' in 2020 IEEE Symposium Series on Computational Intelligence, pp. 1852-1858, doi: 10.1109/SSCI47803.2020.\\
9308140. \href{https://ieeexplore.ieee.org/document/9308140}{\nolinkurl{URL}}

\item Melnikov, T., Lazarev, E., Berezovskaya, O., Chuprov, S., \& Viksnin, I. (2020). ``Empirical Study on Premises Monitoring Algorithm Implementation in Mobile Robotic System'' in The International Conference ``Nonlinearity, Information and Robotics'', pp. 1-6, doi: 10.1109/NIR50484.2020.9290188. \href{https://ieeexplore.ieee.org/document/9290188}{\nolinkurl{URL}}

\item Usova, M., Chuprov, S., \& Viksnin, I. (2020). ``Informational Space and Messages Interaction Models for Smart Factory Concept'' in 2020 IEEE International Workshop on Metrology for Industry 4.0 \& IoT, pp. 617-621, doi: 10.1109/MetroInd4.0IoT48571.2020.9138292. \href{https://ieeexplore.ieee.org/abstract/document/9138292/}{\nolinkurl{URL}}

\item Chuprov, S.,  Marinenkov, E., Viksnin, I., Reznik, L., \& Khokhlov, I (2020). ``Image Processing in Autonomous Vehicle Model Positioning and Movement Control'' in IEEE 6th World Forum on Internet of Things (WF-IoT) Proceedings, pp. 1-6, doi: 10.1109/WF-IoT48130.2020.9221258. \href{https://ieeexplore.ieee.org/document/9221258}{\nolinkurl{URL}}

\item Chuprov, S., Viksnin, I., Kim, I., Reznik, L., \& Khokhlov, I. (2020). ``Reputation and Trust Models with Data Quality Metrics for Improving Autonomous Vehicles Traffic Security and Safety'' in Proc. IEEE/NDIA/INCOSE Syst. Secur. Symp, pp. 1-8, doi: 10.1109/SSS47320.2020.9174269. \href{https://ieeexplore.ieee.org/document/9174269}{\nolinkurl{URL}}

\item Chuprov, S., Viksnin, I., \& Kim, I. (2019) ``Urban Intersection Management with Connected Infrastructure Objects and Autonomous Vehicles'' in 2019 IEEE International Conference on Connected Vehicles and Expo (ICCVE), pp. 1-4, doi: 10.1109/ICCVE45908.201\\
9.8964917. \href{https://ieeexplore.ieee.org/abstract/document/8964917}{\nolinkurl{URL}}

\item Chuprov, S., Viksnin, I., Kim, I., \& Nedosekin, G. (2019). ``Optimization of Autonomous Vehicles Movement in Urban Intersection Management System'' in 2019 24th Conference of Open Innovations Association (FRUCT), pp. 60-66, doi: 10.23919/FRUCT.2019.8711967. \href{https://ieeexplore.ieee.org/abstract/document/8711967/}{\nolinkurl{URL}}

\item Chuprov, S., Viksnin, I., Kim, I., \& Usova, M. (2019). ``Intersection Management Tasks in Mobile Robotic System with Decentralized Control'' in CEUR Workshop Proceedings, vol. 2344, pp. 1-12. \href{http://ceur-ws.org/Vol-2344/paper7.pdf}{\nolinkurl{URL}}

\item Usova, M., Chuprov, S., Viksnin, I., Gataullin, R., Komarova, A., \& Iuganson, A. (2019). ``Model of Smart Manufacturing System'' in International Symposium on Intelligent and Distributed Computing, pp. 356-362, doi: 10.1007/978-3-030-32258-8\_42. \href{https://link.springer.com/chapter/10.1007/978-3-030-32258-8_42}{\nolinkurl{URL}}

\item Marinenkov, E., Chuprov, S., Viksnin, I., \& Kim, I. (2019). ``Empirical Study on Trust, Reputation, and Game Theory Approach to Secure Communication in a Group of Unmanned Vehicles'' in CEUR Workshop Proceedings, vol. 2590, pp. 1-12. \href{http://ceur-ws.org/Vol-2590/paper28.pdf}{\nolinkurl{URL}}

\item Usova, M., Chuprov, S., Viksnin, I., \& Baranova, O. (2019). ``Model of Secure Informational Messages for Ensuring Informational Interaction in Smart Factory'' in CEUR Workshop Proceedings, vol. 2590, pp. 1-8. \href{http://ceur-ws.org/Vol-2590/short28.pdf}{\nolinkurl{URL}}

\item Viksnin, I., Lyakhovenko, J., Tursukov, N., Chuprov, S., \& Sozinova, E. (2019). ``Empirical Study on Modeling of People Behavior in Emergency'' in CEUR Workshop Proceedings, vol. 2590, pp. 1-8. \href{http://ceur-ws.org/Vol-2590/short33.pdf}{\nolinkurl{URL}}

\item Kim, I., Matos-Carvalho, J. P., Viksnin, I., Campos, L. M., Fonseca, J. M., Mora, A., \& Chuprov, S. (2019). ``Use of Particle Swarm Optimization in Terrain Classification based on UAV Downwash'' in 2019 IEEE Congress on Evolutionary Computation (CEC), pp. 604-610, doi: 10.1109/CEC.2019.8790031. \href{https://ieeexplore.ieee.org/abstract/document/8790031/}{\nolinkurl{URL}}

\item Viksnin, I., Chuprov, S., Usova, M., \& Zakoldaev, D. (2019). ``Police Office Model for Multi-agent Robotic Systems'' in IOP Conference Series: Materials Science and Engineering, vol. 497, no. 1, p. 012036, doi: 10.1088/1757-899X/497/1/012036. \href{https://iopscience.iop.org/article/10.1088/1757-899X/497/1/012036/meta}{\nolinkurl{URL}}

%\textbf{Presented and accepted for publication, will be published soon:}
\end{enumerate}
\end{rSection}

\begin{rSection}{Other Publications and Presentations}
\begin{enumerate}
\itemsep-0.5em

\item Chuprov, S. (2020). ``Intersection Management Tasks Application in Mobile Robotic System with Autonomous Vehicle Models'' in Saint-Petersburg's Government Grant Award Winners, pp. 68-74. \href{https://drive.google.com/file/d/1pG5H6agnzNzFAbElviIHcb7Bj3mRQbco/view?usp=sharing}{\nolinkurl{URL}}

\item Chuprov, S. (2020). ``Reputation, Trust, and Data Quality Concept for Security and Safety Assurance in a Group of Autonomous Vehicles'' in IX Young Scientists Congress Proceedings, Electronic Edition. \href{https://kmu.itmo.ru/digests/article/3817}{\nolinkurl{URL}}

\item Chuprov, S. (2019). ``Security Assurance of Mobile Robotic Systems with Decentralized Control'' in 2019 ITMO University Annotated Master's Research Thesis Proceedings, pp. 66-72. \href{https://drive.google.com/file/d/1N9RQR0rh8hOYZR2xX7PoRvpitzYpxaru/view?usp=sharing}{\nolinkurl{URL}}

\item Chuprov, S., Viksnin, I., \& Kim, I. (2019). ``Organization of Safe Intersections Passage by Unmanned Vehicles in a Mobile Robotic System'' in VIII Young Scientists Congress Proceedings, Electronic Edition. \href{https://kmu.itmo.ru/digests/article/209}{\nolinkurl{URL}}

\item Chuprov, S., Viksnin, I., Kim, I., \& Usova, M. (2019). ``Secure Intersection Management Passage with a Group of Autonomous Vehicles in Mobile Robotic System with Decentralized Control'' in ITMO University Young Scientists Almanac, vol. 2, pp. 43-48. \href{https://research.itmo.ru/file/stat/194/almanah_tom_2_2019_(1).pdf}{\nolinkurl{URL}}

\item Chuprov, S., Viksnin, I., \& Usova, M. (2018). ``Police Office Model with Partial Decentralized Control in Mobile Robotic System'' in Regional Informatics and Information Secutity SPOISU Proceedings, vol. 6, pp. 245-249. \href{http://spoisu.ru/files/riib/riib_6_2018.pdf}{\nolinkurl{URL}}
\end{enumerate}

\textbf{Presentations, forums, workshops, and seminars without publications:}
\begin{enumerate}
\itemsep-0.5em

\item Chuprov, S. (2022). ``Discovering and Addressing Privacy and Robustness Flaws in Federated Learning'' at the Research Idea Ring (RIR) in Computer Science, 17 November 2022, Rochester Institute of Technology, Rochester, NY, USA.

\item Chuprov, S., \& Reznik, L. (2021). ``Reputation and Trust Models with Data Quality Metrics for Improving Autonomous Vehicles Traffic Security and Safety'' at Great Lakes Security Day (GLSD) 2021 Online, 12 November 2021, Rochester, NY, USA. -- \textit{Poster session.}

\item Chuprov, S. (2020). ``Security and Safety Approaches to “Soft” Attacks Detection in Intelligent Systems and IoT'' at Technology Development Trend of Intelligent Internet online workshop, 4 December 2020, Saint-Petersburg, Russia.  

\item Chuprov, S. (2020). ``Reputation and Trust Approach for Cyber-physical Systems Security and Safety Assurance'' at ITMO University's XLIX Research, Educational and Methodical Conference, 29 January 2020 - 1 Febeuary 2020. Saint-Petersburg, Russia. 

\item Chuprov, S. (2019). ``Secure and Optimal Intersections' Traversal by a Group of Autonomous Vehicles in Mobile Robotic System'' at Saint-Petersburg's Academic and Industrial Institutions Students' Grant Award Winners Round Table Discussion ``Natural and Exact Sciences'', 6th December 2019. Saint-Petersburg, Russia.    

\item Chuprov, S. (2019). ``Reputation and Trust-based Approach for Security and Safety Assurance in a Group of Autonomous Vehicles'' at Huawei Trust Worthy Software Technology Forum, 28 - 29 November 2019. Saint-Petersburg, Russia. 

\item Chuprov, S. (2019). ``Multi-Agent Physical Testbed for Smart City's Control and Security Algorithms Verification'' at Geek Picnic, 11 - 12 July 2019. Saint-Petersburg, Russia.

\item Chuprov, S. (2018). ``Multi-agent Systems Laboratory Current Research Projects'' at Saint-Petersburg's International Youth Forum 6.0, 1 December 2018. Saint-Petersburg, Russia. 
\end{enumerate}
\end{rSection}


%%%
\begin{comment}
\begin{rSection}{Summary}
Second-year PhD student in Computer Science. Interesting in Cyber-physical systems and Smart City applications security and safety research, working on the doctoral thesis: Models and Methods for Security and Safety Assurance in Distributed Mobile Robotic Systems.  
I want to advance my research career through international cooperation and PhD study in Computer Science/Cybersecurity.  
\end{rSection}
\end{comment}

%\begin{rSection}{Personal Traits}
%\item Great time and project management skills. Like to allocate time in a manner to always keep up. 
%\item Strong teamwork \& collaboration skills. Able to work as an individual as well as in group.
%\item Ability to build relationships and quickly adapt to a new team.
%\item Fast learner. Like to analyze and penetrate the essence of the processes.
%\end{rSection}
\end{document}
